% Generated draft: requires mathematical and source review.
\documentclass[11pt]{article}
\usepackage[margin=25mm]{geometry}
\usepackage{fontspec}
\setmainfont{Latin Modern Roman}
\usepackage{amsmath,amssymb,mathtools,unicode-math}
\setmathfont{Latin Modern Math}
\usepackage{enumitem,needspace,xurl,hyperref}
\setlength{\parindent}{0pt}
\setlength{\parskip}{5pt}
\setlength{\emergencystretch}{3em}
\newcommand{\N}{\mathbb{N}}
\newcommand{\Z}{\mathbb{Z}}
\newcommand{\Q}{\mathbb{Q}}
\newcommand{\R}{\mathbb{R}}
\newcommand{\C}{\mathbb{C}}
\newcommand{\F}{\mathbb{F}}
\newcommand{\A}{\mathbb{A}}
\newcommand{\PP}{\mathbb{P}}
\newcommand{\E}{\mathbb{E}}
\newcommand{\eps}{\varepsilon}
\newcommand{\dd}{\,\mathrm{d}}
\begin{document}
\section*{UnsolvedMath ID: 10900059}
\textbf{Source problem code:} catalogue entry
\subsection{Definitions and mathematical statement}
In the usual compact orientable setting, a Heegaard splitting of \(M\) is a decomposition \(M=C_1\cup_\Sigma C_2\) into compression bodies with common positive boundary \(\Sigma\); for closed \(M\), the pieces are handlebodies. A compressing disk is a properly embedded disk in one side with boundary essential in \(\Sigma\). A splitting is strongly irreducible when each side has a compressing disk and every pair of opposite-side compressing disks \(D_1,D_2\) satisfies \(\partial D_1\cap\partial D_2\ne\varnothing\).

Schleimer asks for a classification of the strongly irreducible Heegaard splittings of a given \(3\)-manifold. In the isotopy interpretation supported by the summary, the target is an explicit description of
\[
 \{\Sigma\subset M:\Sigma\text{ is a strongly irreducible Heegaard surface}\}
 \big/\text{ambient isotopy in }M,
\]
with complete criteria distinguishing the represented classes. No genus bound or assertion that this set is finite is supplied. Whether the requested classification is effective/algorithmic, and whether splittings retain an ordering of the two sides, is not stated.

The source links arXiv:1509.05945 as progress and reports that its abstract solves the classification problem for non-Haken hyperbolic \(3\)-manifolds. Its summary describes finiteness-related results there and says the general Haken case is not resolved. This restricted reported result is not promoted to a classification for all \(3\)-manifolds. The original statement does not specify compactness, orientability, boundary conventions or low-genus conventions, so the definitions above record the standard context rather than resolve those omissions.
\subsection{Short English statement}
Describe all strongly irreducible Heegaard splittings of a fixed \(3\)-manifold up to the intended equivalence, keeping the reported non-Haken hyperbolic result separate from the general question.
\subsection{Sources}
\begin{itemize}
\item Schleimer, attributed item 8.5 in Delp, Hoffoss and Manning, \emph{Problems In Groups, Geometry, and Three-Manifolds} (2015), supplied PDF page 6. \url{https://arxiv.org/pdf/1512.04620#page=6}; original-status link \url{https://arxiv.org/abs/1512.04620}. Recorded access 2026-07-29; extraction source-tex checked against rendered PDF text, as reported in the record.
\item Source-linked progress: \url{https://arxiv.org/abs/1509.05945}; the record reports its abstract as resolving the non-Haken hyperbolic case, but supplies no complete author/title citation. The summary mentions Namazi--Ishikawa and other deformation-space work without precise references.
\item Local UnsolvedMath registry, \texttt{data/unsolvedmath/problems.json}, permanent ID \(10900059\), exact source code \texttt{catalogue entry}; supplied partially\_solved/PARTIAL-PROGRESS status.
\end{itemize}
\end{document}
